% ===============================================================
% AB-05b: Reguläre Verben auf -er/-ir
% Abgeleitet aus LE-05b (Score 9.3/10)
% ===============================================================

\documentclass[11pt, a4paper]{article}

% ===============================================================
% SW-MODUS TOGGLE
% ===============================================================
\newif\ifswmodus
\swmodusfalse    % Standard: Farbe

\usepackage[utf8]{inputenc}
\usepackage[T1]{fontenc}
\usepackage[ngerman]{babel}
\usepackage{helvet}
\renewcommand{\familydefault}{\sfdefault}

\usepackage[margin=2cm]{geometry}
\usepackage{enumitem}
\usepackage{tabularx}
\usepackage{booktabs}

\usepackage{xcolor}
\usepackage{tcolorbox}
\tcbuselibrary{skins}

\usepackage{fancyhdr}
\usepackage{lastpage}
\usepackage{needspace}

% Farben - Basis
\definecolor{spanisch-color}{HTML}{c41e3a}
\definecolor{grau-color}{HTML}{888888}
\definecolor{hellgrau-color}{HTML}{f5f5f5}
\definecolor{orange-color}{HTML}{b35c00}
\definecolor{blau-color}{HTML}{2c5aa0}
\definecolor{gruen-color}{HTML}{2a7a4b}

% SW-Farben
\definecolor{sw-dark}{HTML}{000000}
\definecolor{sw-gray}{HTML}{666666}
\definecolor{sw-white}{HTML}{ffffff}

% Bedingte Farbzuweisung
\ifswmodus
    \colorlet{spanisch}{sw-dark}
    \colorlet{grau}{sw-gray}
    \colorlet{hellgrau}{sw-white}
    \colorlet{orange}{sw-dark}
    \colorlet{blau}{sw-gray}
    \colorlet{gruen}{sw-gray}
\else
    \colorlet{spanisch}{spanisch-color}
    \colorlet{grau}{grau-color}
    \colorlet{hellgrau}{hellgrau-color}
    \colorlet{orange}{orange-color}
    \colorlet{blau}{blau-color}
    \colorlet{gruen}{gruen-color}
\fi

\newcommand{\fachfarbe}{spanisch}

% Variablen
\newcommand{\thema}{Reguläre Verben auf -er/-ir}
\newcommand{\sequenz}{05}
\newcommand{\block}{b}
\newcommand{\fach}{Spanisch}
\newcommand{\klasse}{7}
\newcommand{\maxpunkte}{20}

% Kopf-/Fußzeile
\pagestyle{fancy}
\fancyhf{}
\renewcommand{\headrulewidth}{0.4pt}
\renewcommand{\footrulewidth}{0.4pt}
\lhead{\fach{} -- Klasse \klasse}
\chead{AB-\sequenz\block}
\rhead{}
\lfoot{}
\rfoot{Seite \thepage\ von \pageref{LastPage}}

% Befehle
\newcommand{\punktebox}[1]{\hfill\fbox{\small \_/#1}}
\newcommand{\luecke}[1]{\rule{#1}{0.4pt}}
\newcommand{\es}[1]{\textcolor{\fachfarbe}{\textbf{#1}}}
\newcommand{\er}[1]{\textcolor{blau}{\textbf{#1}}}
\newcommand{\ir}[1]{\textcolor{gruen}{\textbf{#1}}}

% Merkkasten
\newtcolorbox{merkkasten}[1][]{
    colback=hellgrau,
    colframe=\fachfarbe,
    colbacktitle=white,
    coltitle=\fachfarbe,
    boxrule=1pt,
    arc=0pt,
    left=8pt, right=8pt, top=8pt, bottom=8pt,
    fonttitle=\bfseries,
    attach boxed title to top left={yshift=-2mm, xshift=5mm},
    boxed title style={boxrule=0pt, colback=white},
    title=#1
}

% Hinweis
\newtcolorbox{hinweis}{
    colback=white,
    colframe=orange,
    colbacktitle=white,
    coltitle=orange,
    boxrule=1pt,
    arc=0pt,
    left=5pt, right=5pt, top=5pt, bottom=5pt,
    fonttitle=\bfseries,
    attach boxed title to top left={yshift=-2mm, xshift=5mm},
    boxed title style={boxrule=0pt, colback=white},
    title=Tipp
}

% Aufgabe
\newenvironment{aufgabe}[1]{%
    \needspace{5\baselineskip}%
    \nopagebreak[4]%
    \vspace{10pt}
    \noindent\textbf{\textcolor{\fachfarbe}{Aufgabe #1}}
    \nopagebreak[4]%
    \vspace{5pt}
    \nopagebreak[4]%
    \begin{list}{}{\leftmargin=0pt}
    \item[]
}{%
    \end{list}
}

% ===============================================================
\begin{document}

% Kopfbereich
\begin{center}
    {\Large\bfseries\textcolor{\fachfarbe}{Arbeitsblatt: \thema}}
\end{center}

\vspace{5pt}

\noindent
\begin{tabularx}{\textwidth}{|X|X|X|}
\hline
\textbf{Name:} \luecke{4cm} &
\textbf{Klasse:} \klasse &
\textbf{Datum:} \luecke{2cm} \\
\hline
\end{tabularx}

\vspace{15pt}

% ===============================================================
% MERKKASTEN 1: Verben auf -er
% ===============================================================

\begin{merkkasten}[Merkkasten 1: Verben auf \er{-er} (z.B. comer{,} beber{,} leer{,} correr)]

\begin{center}
\begin{tabular}{|l|l|l|}
\hline
\textbf{Person} & \textbf{Pronomen} & \textbf{comer} (essen) \\
\hline
1. Sg. & yo & \luecke{2.5cm} \\
\hline
2. Sg. & tú & \luecke{2.5cm} \\
\hline
3. Sg. & él/ella & \luecke{2.5cm} \\
\hline
1. Pl. & nosotros/-as & \luecke{2.5cm} \\
\hline
2. Pl. & vosotros/-as & \luecke{2.5cm} \\
\hline
3. Pl. & ellos/ellas & \luecke{2.5cm} \\
\hline
\end{tabular}
\end{center}

\vspace{0.3em}
\textbf{Endungen -er:} \er{-o, -es, -e, -emos, -éis, -en}
\end{merkkasten}

\vspace{10pt}

% ===============================================================
% MERKKASTEN 2: Verben auf -ir
% ===============================================================

\begin{merkkasten}[Merkkasten 2: Verben auf \ir{-ir} (z.B. vivir{,} escribir{,} abrir)]

\begin{center}
\begin{tabular}{|l|l|l|}
\hline
\textbf{Person} & \textbf{Pronomen} & \textbf{vivir} (leben/wohnen) \\
\hline
1. Sg. & yo & \luecke{2.5cm} \\
\hline
2. Sg. & tú & \luecke{2.5cm} \\
\hline
3. Sg. & él/ella & \luecke{2.5cm} \\
\hline
1. Pl. & nosotros/-as & \luecke{2.5cm} \\
\hline
2. Pl. & vosotros/-as & \luecke{2.5cm} \\
\hline
3. Pl. & ellos/ellas & \luecke{2.5cm} \\
\hline
\end{tabular}
\end{center}

\vspace{0.3em}
\textbf{Endungen -ir:} \ir{-o, -es, -e, -imos, -ís, -en}
\end{merkkasten}

\vspace{10pt}

\begin{hinweis}
\textbf{Unterschied -er und -ir:} Nur bei \textbf{nosotros} und \textbf{vosotros}!\\
-er: com\er{emos}, com\er{éis} \hspace{1cm} -ir: viv\ir{imos}, viv\ir{ís}
\end{hinweis}

\vspace{15pt}

% ===============================================================
% AUFGABE 1: Konjugation -er (AFB I/II)
% ===============================================================

\begin{aufgabe}{1}
Setze die richtige Form des Verbs in Klammern ein. \punktebox{5}
\end{aufgabe}

\begin{enumerate}[label=\alph*)]
    \item Yo \luecke{2.5cm} pizza. (comer)
    \item Tú \luecke{2.5cm} agua. (beber)
    \item María \luecke{2.5cm} un libro. (leer)
    \item Nosotros \luecke{2.5cm} en el parque. (correr)
    \item Ellos \luecke{2.5cm} mucha leche. (beber)
\end{enumerate}

\vspace{15pt}

% ===============================================================
% AUFGABE 2: Konjugation -ir (AFB I/II)
% ===============================================================

\begin{aufgabe}{2}
Setze die richtige Form des Verbs in Klammern ein. \punktebox{5}
\end{aufgabe}

\begin{enumerate}[label=\alph*)]
    \item Yo \luecke{2.5cm} en Rostock. (vivir)
    \item Tú \luecke{2.5cm} un mensaje. (escribir)
    \item Pedro \luecke{2.5cm} la ventana. (abrir)
    \item Nosotros \luecke{2.5cm} en Alemania. (vivir)
    \item Vosotros \luecke{2.5cm} cartas. (escribir)
\end{enumerate}

\vspace{15pt}

% ===============================================================
% AUFGABE 3: Sätze bilden (AFB II)
% ===============================================================

\begin{aufgabe}{3}
Bilde Sätze mit den angegebenen Elementen. Konjugiere das Verb! \punktebox{4}
\end{aufgabe}

\begin{hinweis}
\textbf{Vokabeln:} en el parque = im Park, un libro = ein Buch, en casa = zu Hause, con amigos = mit Freunden
\end{hinweis}

\vspace{0.5em}

\begin{enumerate}[label=\alph*)]
    \item yo / leer / un libro / en casa\\[0.8em]
    \luecke{12cm}

    \item mi hermano / correr / en el parque\\[0.8em]
    \luecke{12cm}

    \item nosotros / comer / pizza / con amigos\\[0.8em]
    \luecke{12cm}

    \item tú / vivir / en España\\[0.8em]
    \luecke{12cm}
\end{enumerate}

\vspace{15pt}

% ===============================================================
% AUFGABE 4: Vergleich -ar/-er/-ir (AFB II/III)
% ===============================================================

\begin{aufgabe}{4}
Vervollständige die Tabelle und erkläre den Unterschied. \punktebox{6}
\end{aufgabe}

\begin{center}
\begin{tabular}{|l|c|c|c|}
\hline
\textbf{Person} & \textbf{hablar} (-ar) & \textbf{comer} (-er) & \textbf{vivir} (-ir) \\
\hline
yo & hablo & \luecke{1.8cm} & \luecke{1.8cm} \\
\hline
tú & hablas & \luecke{1.8cm} & \luecke{1.8cm} \\
\hline
él/ella & habla & \luecke{1.8cm} & \luecke{1.8cm} \\
\hline
nosotros & hablamos & \luecke{1.8cm} & \luecke{1.8cm} \\
\hline
vosotros & habláis & \luecke{1.8cm} & \luecke{1.8cm} \\
\hline
ellos & hablan & \luecke{1.8cm} & \luecke{1.8cm} \\
\hline
\end{tabular}
\end{center}

\vspace{1em}

\textbf{Erkläre:} Wo unterscheiden sich -er und -ir? (2P)

\noindent\rule{\textwidth}{0.4pt}\vspace{8pt}
\noindent\rule{\textwidth}{0.4pt}

\vspace{20pt}
\hrule
\vspace{10pt}

\begin{flushright}
\textbf{Punkte:} \luecke{1cm} / \maxpunkte
\end{flushright}

\end{document}
